\documentclass[12pt,a4paper]{article}
\usepackage[utf8]{inputenc}
\usepackage[T2A]{fontenc}
\usepackage[russian]{babel}
\usepackage{amsmath,amssymb,amsthm}
\usepackage{booktabs,array}
\usepackage{geometry}
\usepackage{microtype}
\geometry{margin=2.3cm}

\newtheorem{protocol}{Протокол}
\newtheorem{nogocriterion}{Критерий провала}

\title{Коллективный $\pi$-atlas gate:\\
общая база без доказанной уникальности}
\author{}
\date{5 августа 2026 г.}

\begin{document}
\maketitle

\section{Ретроспективный тест}

Одиннадцать уже объявленных атласных формул переписаны с общей переменной
$x$ вместо $\pi$. В тест включены $\alpha_s$, $\sin^2\theta_W$, четыре
массовых отношения, $V_{cb}$ и три космологических доли. Коэффициенты и
назначения формул не оптимизируются; меняется только общий базис
$x\in[2,4]$.

Метрика равна
\begin{equation}
 E(x)=\sqrt{\frac1{11}\sum_{i=1}^{11}
 \left[\log\frac{F_i(x)}{O_i}\right]^2}.
\end{equation}
Получено
\begin{equation}
 E(\pi)=0.01827893,
 \qquad
 x_{best}=3.16422,
 \qquad
 E(x_{best})=0.01587965.
\end{equation}
На равномерной сетке $[2,4]$ значение $\pi$ входит в лучшие $2.27\%$ баз.
Это заметно сильнее одиночного совпадения, но не является уникальностью:
\begin{equation}
 E(\sqrt{10})=0.01589845<E(\pi),
 \qquad
 E(22/7)\approx0.01803<E(\pi).
\end{equation}

Leave-one-out проверка сохраняет оптимум в узком интервале
\begin{equation}
 3.15535\le x_{best}^{(-i)}\le3.17735.
\end{equation}
Следовательно, эффект не создаётся одной формулой. Он указывает на устойчивую
общую арифметическую базу около $3.16$, но данные не выделяют именно $\pi$.

\begin{nogocriterion}[Запрет ретроспективного blind-статуса]
Формулы и наблюдаемые были собраны после знакомства с численными значениями.
Поэтому процентиль $\pi$ измеряет внутреннюю компрессию существующего атласа,
а не вероятность физической гипотезы. Он не повышает формулы до предсказаний.
\end{nogocriterion}

\section{Замороженный prospective-протокол}

\begin{protocol}[Условия следующего теста]
Для будущего blind-аудита фиксируются: одиннадцать текущих формул, равные веса,
RMS логарифмической ошибки и интервал $x\in[2,4]$. После раскрытия новой
наблюдаемой запрещены изменение коэффициентов, переназначение формулы,
удаление провала и добавление correction-term. Главный блокер состоит в том,
что назначение новой наблюдаемой формуле должно быть выведено операторным
selector-ом до знания её значения.
\end{protocol}

Таким образом, коллективный атлас сохраняет исследовательский интерес: общая
база около $3.16$ действительно сжимает несколько секторов одновременно.
Однако следующий шаг --- не добавление двенадцатой красивой формулы, а вывод
правила адресации observables к операторам.

Полный протокол сохранён в
\texttt{s2t\_collective\_pi\_atlas\_base\_audit.py} и
\texttt{s2t\_collective\_pi\_atlas\_base\_results.json}.

\end{document}